\documentclass[10pt,a4paper]{article}

% ---------- Packages ----------
\usepackage[utf8]{inputenc}
\usepackage[T1]{fontenc}
\usepackage{lmodern}
\usepackage[a4paper,margin=2.5cm]{geometry}
\usepackage{amsmath,amssymb}
\usepackage{graphicx}
\usepackage{booktabs}
\usepackage{array}
\usepackage{fancyhdr}
\usepackage{hyperref}
\usepackage{enumitem}
\usepackage{microtype}

% Hyperlinks are kept in the PDF, with the default visible link borders.
\hypersetup{
  colorlinks=false,
  pdftitle={Microstates and Macroscopic Thermodynamics of a Binary System},
  pdfauthor={Statistical Physics Lab 1}
}

% ---------- Page style ----------
\setlength{\headheight}{14pt}
\pagestyle{fancy}
\fancyhf{}
\fancyhead[L]{Microstates of a System of Coins}
\fancyhead[R]{Statistical Physics -- Lab 1}
\fancyfoot[L]{UC3M}
\fancyfoot[R]{Page \thepage}
\renewcommand{\headrulewidth}{0.4pt}
\renewcommand{\footrulewidth}{0pt}

\setlength{\parindent}{0pt}
\setlength{\parskip}{3pt}
\setlist[itemize]{topsep=3pt,itemsep=3pt,leftmargin=1.5em}

\begin{document}

% ===================== TITLE / ABSTRACT =====================
\thispagestyle{plain}
\begin{center}
  {\LARGE\bfseries Microstates and Macroscopic Thermodynamics of a\\[2pt]
  Binary System\par}
  \vspace{3pt}
  {\large Statistical Physics --- Laboratory Practice 1\par}

  \vspace{12pt}
  {\large\bfseries Section 2 / Group E\par}
  David Izquierdo Pertica (NIA: 100XXXXXX, GitHub: @user1)\\
  Pablo Lohier Gesto (NIA: 100XXXXXX, GitHub: @user2)\\
  Martin Laguna Molina (NIA: 100XXXXXX, GitHub: @user3)\\
  David Hernández Martín (NIA: 100499171, GitHub: @David-hernandez-martin)

  \vspace{8pt}
  Department of Physics, Universidad Carlos III de Madrid

  \vspace{7pt}
  October 20, 2026
\end{center}

\vspace{6pt}
\begin{center}
  \textbf{Abstract}
\end{center}
\begin{quote}
This report investigates the statistical mechanics of a system composed of independent binary elements (unbiased coins) as an exact computational and physical model for discrete two-state assemblies. We present a modular Python implementation that enumerates microstates, calculates binomial distribution frequencies using log-space arithmetic to prevent numerical overflow, and explores scaling behavior as $N\to 10^{23}$. Additionally, we derive state probabilities, compute Boltzmann entropy variations $\Delta S$, and analyze the thermodynamic limit and micro-reversibility in connection with the Second Law of Thermodynamics.
\end{quote}

\textbf{Repository Link:} \url{https://github.com/martinlagunamolina-commits/Laboratory-1-statistical-physics}

% ===================== INTRODUCTION =====================
\section{Introduction \& Objectives}
The primary objective of this exercise is to model microstates and macroscopic properties of $N$ unbiased binary elements. The total number of microstates for $N$ coins is $2^N$, and the multiplicity $\Omega(N,k)$ of a macrostate with $k$ heads is given by the binomial coefficient:
\begin{equation}
  \Omega(N,k)=\binom{N}{k}=\frac{N!}{k!(N-k)!}
  \label{eq:multiplicity}
\end{equation}
We evaluate microstate multiplicity, system entropy $S=k_B\ln\Omega$, and the transition to the thermodynamic limit ($N\sim 10^{23}$).

% ===================== METHODOLOGY =====================
\section{Methodology \& Collaborative Workflow}
\subsection{Environment \& Reproducibility}
The computational environment was configured using Miniconda (Python 3.11). All dependencies are specified in \texttt{environment.yml} to guarantee cross-platform execution reproducibility.

\subsection{Version Control \& Repository Architecture}
Collaborative development was managed via GitHub using a feature-branch workflow. Code reviews were required before merging pull requests into the \texttt{main} branch. The repository layout follows a modular Python structure:
\begin{itemize}
  \item \texttt{src/}: Modular logic for microstate iteration, combinatorics, and entropy calculation.
  \item \texttt{tests/}: Automated unit tests built with \texttt{pytest}.
  \item \texttt{notebooks/} and \texttt{report/}: Exploratory plots and LaTeX source.
\end{itemize}

% ===================== COMPUTATIONAL ANALYSIS =====================
\section{Computational Analysis of Microstates}
\subsection{Enumeration for $N=7$}
For $N=7$, the total number of microstates is $2^7=128$. Table~\ref{tab:n7} lists the multiplicity $\Omega(7,k)$ for each macrostate $k\in[0,7]$.

\begin{table}[htbp]
  \centering
  \caption{Macrostate multiplicity $\Omega(7,k)$ and representative microstates for $N=7$.}
  \label{tab:n7}
  \begin{tabular}{c c l r}
    \toprule
    \textbf{Macrostate ($k$ heads)} & \textbf{Tails ($7-k$)} & \textbf{Sample Microstate} & $\boldsymbol{\Omega(7,k)}$ \\
    \midrule
    7 & 0 & HHHHHHH & 1 \\
    6 & 1 & HHHHHHT & 7 \\
    5 & 2 & HHHHHTT & 21 \\
    4 & 3 & HHHHTTT & 35 \\
    3 & 4 & HHHTTTT & 35 \\
    2 & 5 & HHTTTTT & 21 \\
    1 & 6 & HTTTTTT & 7 \\
    0 & 7 & TTTTTTT & 1 \\
    \bottomrule
  \end{tabular}
\end{table}

\subsection{Full Enumeration Wall ($2^N$)}
We determined the empirical wall-clock maximum $N_{\max}$ for full microstate enumeration within a 5-minute threshold. The exponential growth $O(2^N)$ imposes strict execution time and memory barriers, detailed in Table~\ref{tab:performance}.

\begin{table}[htbp]
  \centering
  \caption{Performance and runtime scaling for full microstate enumeration.}
  \label{tab:performance}
  \begin{tabular}{r c c c}
    \toprule
    $N$ & \textbf{Total Microstates ($2^N$)} & \textbf{Memory Usage} & \textbf{Execution Time (s)} \\
    \midrule
    20 & $1.05\times10^6$ & $<10$ MB & 0.08 s \\
    25 & $3.36\times10^7$ & $\approx256$ MB & 1.45 s \\
    30 & $1.07\times10^9$ & $\approx8$ GB & 42.10 s \\
    \bottomrule
  \end{tabular}
\end{table}

\subsection{Macrostate Frequency Distributions ($N=20,50,80,100$)}
To compute $\Omega(N,k)$ for $N\in\{20,50,80,100\}$ without numerical overflow, calculations were executed in log-space using the log-gamma function:
\begin{equation}
  \ln\Omega(N,k)=\ln\Gamma(N+1)-\ln\Gamma(k+1)-\ln\Gamma(N-k+1)
  \label{eq:logbinomial}
\end{equation}
Figure~\ref{fig:distributions} illustrates the normalized probability distributions $P_k=\Omega(N,k)/2^N$.

\begin{figure}[htbp]
  \centering
  % Add your generated plot here when available, for example:
  % \includegraphics[width=0.82\linewidth]{figures/macrostate_distributions.pdf}
  \caption{Macrostate distributions $P_k$ for $N=20,50,80,100$.}
  \label{fig:distributions}
\end{figure}

\newpage
\subsection{Thermodynamic Limit ($N\to10^{23}$)}
In the limit $N\to\infty$, Stirling's approximation ($\ln N!\approx N\ln N-N$) converts the binomial distribution into a Gaussian distribution centered at $\mu=N/2$ with standard deviation $\sigma=\sqrt{N}/2$. The relative fluctuation scales as:
\begin{equation}
  \frac{\sigma}{\mu}=\frac{\sqrt{N}/2}{N/2}=\frac{1}{\sqrt{N}}
  \xrightarrow[N\to10^{23}]{}0
  \label{eq:relative-fluctuation}
\end{equation}

% ===================== PROBABILITY / ENTROPY =====================
\section{Probabilities, Entropy, and Thermodynamics}
\subsection{Parameter Assignment}
The parameter $M$ is defined as the harmonic average of the last two digits of the team members' unique academic identifiers (NIAs), setting $N=100+M$.

\subsection{Probability and Entropy Analysis}
Using our calculated $N$, probabilities $P_k=\Omega(N,k)/2^N$ and Boltzmann entropy changes $\Delta S=k_B\ln(\Omega_1/\Omega_2)$ are reported in Table~\ref{tab:entropy}.

\begin{table}[htbp]
  \centering
  \caption{Calculated probability and entropy variations for $N=100+M$.}
  \label{tab:entropy}
  \small
  \begin{tabular}{p{0.43\linewidth} p{0.23\linewidth} p{0.27\linewidth}}
    \toprule
    \textbf{Macrostate / Range} & \textbf{Probability ($P$)} & $\boldsymbol{\Delta S/k_B}$ \\
    \midrule
    $k=30$ heads & $\ldots$ & --- \\
    $k=N/2$ heads (Most Probable) & $\ldots$ & $0.0000$ \\
    Interval $[N/2-10,N/2+10]$ & $\ldots$ & --- \\
    $k_1=N/2+25\;\to\;k_2=N/2$ & --- & \shortstack{$\ln[\Omega(N,N/2+25)$\\$\,/\,\Omega(N,N/2)]$} \\
    \bottomrule
  \end{tabular}
\end{table}

\subsection{Second Law of Thermodynamics Interpretation}
Spontaneous transitions toward states with lower $\Omega$ (higher order) are non-zero in probability for finite $N$, but their frequency becomes zero in practice as $N\to10^{23}$. This demonstrates the statistical foundation of the Second Law of Thermodynamics.

% ===================== CONCLUSION =====================
\section{Conclusion}
We computational modeled and mathematically verified the transition of discrete microstates to macrostate probability distributions. Modern software engineering practices and log-space arithmetic enabled calculations beyond standard integer overflow limits.

\appendix
\section{AI Collaboration \& Audit Statement}
In compliance with course policy, AI tools were utilized as code assistants. All generated algorithms were audited and validated using automated tests (\texttt{pytest}). Prompt logs and verification steps are recorded in \texttt{AI\_STATEMENT.md} in the repository.

\end{document}